\subsection{The actual squares for unit terminals of lengths $(2,4)$ and $(2,5)$}\label{proof:Q5-C069}
The following square and moment calculations are used in the reference proofs. Put $r=Kg$, $s=K^2g$, $t_0=K^3g$ and $m=1-d$.
For $(2,4)$, the actual positive paths satisfy
$h_2=q-r$, $f_4=p+r+s+t_0-d-dp$, so
$h_2+f_4=2m+R$, $R=s+t_0-dg$.
The actual product is at most $F_0=(m+R/2)^2$.
For any such square write
\begin{equation}\tag{L69.10.square}\label{eq:L69:square}
\begin{gathered}
Y=TF_0=m^2+w_0+e,\quad w_0=-m^2g+mD_0-mKR,\quad D_0=\E R,\\
e=T(R^2)/4,\quad0\le e\le V/2,\quad V=\E R^2.
\end{gathered}
\end{equation}
For $(2,4)$ the exact original moments are
$\E R=\delta_2$, $\E KR=d^2+\delta_2$,
$\E(rR)=-\delta_3$, $\E(sR)=\delta_3+d\delta_2$.
Consequently
\begin{equation}\tag{L69.10.2}\label{eq:L69:firstfour}
\E[f_2(w_0+e)]=-m^2d-md^2+\E[f_2e]
\le-m^2d-md^2+V/2,
\end{equation}
where the actual adjoint gives
$\E[f_2e]=\E[f_3R^2]/4\le V/2$.
The complete second moment, using $m+d=1$, is
\begin{equation}\tag{L69.10.3}\label{eq:L69:secondfour}
\E w_0^2=m^2[m^2d+2\delta_3+\delta_4+\delta_5
-d^2\delta_2-\delta_2^2].
\end{equation}
On $0<d\le1/4$ set
$v=15/64+(13/16)d+d^2$,
$H=m^2+17/40-(4/25)d-(2/5)d^2$.
The exact $V=d^3+2d\delta_2+\delta_3+\delta_4$ gives
$V\le dv$ and $V\le d(1/2+d)^2$,
using $\theta^2+\theta^3=7605/32768<15/64$.
For the bracket in \eqref{eq:L69:secondfour}, normalized by $d$,
write $v_1=\delta_2/d$. Its upper bound is
$m^2+(2\theta+\theta^2+\theta^3)v_1-d^2v_1-dv_1^2$,
increasing on $[0,\theta]$ since its derivative is at least
$2\theta-d^2-2d\theta>0$.
At $v_1=\theta$ use
$2\theta^2+\theta^3+\theta^4=444977/1048576<17/40$,
$\theta\ge2/5$, $\theta^2\ge4/25$ to obtain
$\E w_0^2\le m^2dH$.
$0<H\le57/40<(6/5)^2$ and Cauchy--Schwarz with $e\le V/2$ give
\begin{equation}\tag{L69.10.4}\label{eq:L69:fourvariance}
\E(w_0+e)^2\le m^2dH+(6/5)md^{3/2}v+d^2v^2/4.
\end{equation}

We supply the full original-$f_2$ curvature certificate.
With $x=\sqrt d\in[0,1/2]$, the row bound gives
$Y\le m^2q+B_0$, $B_0/m\le x/2+(41/32)x^2
\le C=x/2+(21/16)x^2$.
Since $f_2\le p+x=2+x-q$, its fourth Taylor coefficient divided
by $m^2$ is at most
$Q(q)=(2+x-q)\{[m(q+1)+C]^2+2m^2\}$.
Put $c=C/m\le37x/24<1$.
$Q'(0)>0$ and $Q'(1/2)/m^2=1/4+3x+2xc-c^2>0$;
$Q'$ is concave, while
$Q'(1)/m^2=-2+4x+(-2+2x)c-c^2\le0$ and $Q''<0$ after one.
Thus its maximum is in $[1/2,1]$, where $Q''\le-2m^2$.
At $q=3/4$ the normalized derivative is
$D=-11/16+(7/2)x+(2x-1)c-c^2$, with $|D|\le7/10$.
For the lower bound use $c+c^2\le(3145/1152)x<(7/2)x$.
For $x\le2/5$, use $c\ge x/2$ to get
$D\le269/400<7/10$; for $x\ge2/5$, $c/x$ increases and is
at least $205/168>61/50$, giving
$D\le-11/16+(57/25)x+(2379/2500)x^2\le863/1250<7/10$.
Strong concavity gives
\[
\max Q\le\Phi(x)=(5/4+x)
\{[7/4+x/2-7x^2/16]^2+2(1-x^2)^2\}
+(49/400)(1-x^2)^2.
\]
Its coefficients, in ascending degree, are
$10321/1600,\allowbreak 29/4,\allowbreak -621/128-49/200,\allowbreak -373/64,\allowbreak 
2357/1024+49/400,\allowbreak 561/256$.
Absorb $x^4\le x^2/4$, $x^5\le x^3/4$, using
$459827/102400>22/5$, $5407/1024>21/4$, to get
$\Phi\le13/2+(29/4)x-(22/5)x^2-(21/4)x^3$.
This increases up to $2/5$ and there equals $209/25<17/2$.
For $2/5\le x\le1/2$, use $x^3\ge(2/5)x^2$ to get the increasing
upper bound $13/2+(29/4)x-(13/2)x^2$, equal to $17/2$ at $1/2$.
Thus
\begin{equation}\tag{L69.10.5}\label{eq:L69:fourcurvature}
f_2(Y^2+2m^2Y+3m^4)\le(17/2)m^2.
\end{equation}
The exact fourth expansion and the above moments bound the target by
\[
U=m^8(1-5d)-4m^7d^2+2m^6dv+(17/2)m^4dH
+(51/5)m^3d^{3/2}v+(17/8)m^2d^2v^2.
\]
The positive Taylor series gives
$[m^{-4}-m^4(1-5d)]/d\ge13-16d+54d^2+14d^3+61d^4$.
After combining all nonsquare-root terms the normalized deficit is
at least
\[
67/160+(2269/400)d+(6029/160)d^2+(227/8)d^3+55d^4
-(51/5)(\sqrt d/m)v-(17/8)(d/m^2)v^2.
\]
Use $m^{-1}\le1+4d/3$, $m^{-2}\le16/9$.
The first negative cost has $x,x^3,x^5,x^7$ coefficients
$153/64,459/40,85/4,68/5$; the second has $d$ through $d^5$
coefficients $425/2048,\allowbreak 1105/768,\allowbreak 4913/1152,\allowbreak 221/36,\allowbreak 34/9$.
With $d^5\le d^4/4$, $x^5\le x^4/2$, $x^7\le x^6/2$, the
deficit is at least
$67/160-(153/64)x+(27/5)x^2-(459/40)x^3+25x^4$.
The retained quadratic and quartic coefficients are at least
$279807/51200>27/5$, $98371/3840>25$; remaining sixth and eighth
coefficients are nonnegative.
Since $25x^2-(459/40)x+27/5\ge4$, whose square-completion remainder
is $13319/160000>0$, the deficit is at least
$67/160-(153/64)x+4x^2\ge20171/327680$.
Multiplying by $m^4\ge81/256$ proves the low-range bound
$\E[f_2\{T(h_2f_4)\}^4]\le1-d/100$.
The exact actual mean is
$\E(h_2f_4)=1-3d+\delta_2-d^2+\delta_3$.
For $d\ge1/4$ its decreasing upper bound is
$1353/4096<2/5$, so $TF<4/5$ and the target is at most $256/625$.

For $(2,5)$ keep the same $m$ but use the actual retained vector
\begin{equation}\tag{L69.11.vector}\label{eq:L69:fivevector}
R=\sum_{j=3}^5K^j\one-(d+\delta_2)g-\delta_2r.
\end{equation}
Finite positive convolution retaining $f_0,f_1$ and the actual
$\delta_2f_2$ portion gives $h_2+f_5\le2m+R$.
Its multiplier at $g$ is $-A-i\sqrt\xi B$, where
$A=d+\delta_2+\xi-\xi^2$, $B=\delta_2+\xi$.
Let $v=101/800+(13/16)d+(33/16)d^2$ and
$H=m^2+49/128+(3/10)d-(1/2)d^2-(19/20)d^3$.
Here $D_0=(1+d)\delta_2$ and the exact reference cancellation is
\begin{equation}\tag{L69.11.2}\label{eq:L69:fivefirst}
\E[f_2(w_0+e)]=-m^2d-md^2(1+\delta_2)
+m(\delta_3-\delta_4)+\E[f_2e],\quad\E[f_2e]\le V/2.
\end{equation}
The actual difference $\delta_3-\delta_4$ is bounded as one spectral
term: $\xi^2(1-\xi)$ increases to $\theta^2\kappa$, with
$4\theta^2\kappa=3211/8192<2/5$.
The exact vector square is
\[
V=d(d+\delta_2)^2+[2(d+\delta_2)+\delta_2^2]\delta_2
+(1-2d)\delta_3-\delta_4+\delta_5\le dv.
\]
The positive $2\delta_2^2$ cancels with $-2d\delta_3$ using
$\delta_3d\ge\delta_2^2$; the remaining frequency bound gives
$V/d\le\theta^2-\theta^3+\theta^4+2\theta d
+[(1+\theta)^2+\theta^3]d^2$.
Use $131313/1048576<101/800$ and
$66997/32768<33/16$ for the two constants.
The complete second square, writing $u_2=\delta_2$, is
\begin{align*}
\E w_0^2/m^2={}&m^2d+(1-2d-d^2)u_2^2-d^2u_2-2d^3u_2+u_2^3\\
&+[2+(4+2d)u_2+u_2^2]\delta_3+(1-2d)\delta_4-\delta_5+\delta_6.
\end{align*}
The complete high-frequency combination
$2\xi^2+(1-2d)\xi^3-\xi^4+\xi^5$ increases on $[0,\theta]$.
After bounding it, the upper expression increases in $u_2/d$,
with derivative at least $d(4\theta^2-d-2d^2)>0$ on $d\le1/4$.
At $u_2/d=\theta$ it is bounded by
\[
m^2+(2\theta^2+\theta^3-\theta^4+\theta^5)
+\theta^2(1+2\theta)d
-(2\theta^2+\theta-3\theta^3-\theta^4)d^2
-(\theta^2+2\theta)d^3\le H.
\]
The four respective coefficient comparisons are $<49/128$,
$<3/10$, $>1/2$, $>19/20$; the first and third positive differences
have numerators $62403$ over $33554432$ and $8335$ over $1048576$.
$H$ is positive and decreasing, $H\le177/128<(6/5)^2$, hence
\begin{equation}\tag{L69.11.6}\label{eq:L69:fivevariance}
\E(w_0+e)^2\le m^2dH+(6/5)md^{3/2}v+d^2v^2/4.
\end{equation}

The five-arm curvature certificate has, with $x=\sqrt d$,
$\sqrt V\le(3/8+3d/2)\sqrt d$:
the two squared vector constants are
$(\theta-\theta^2)^2+\theta^3=131313/1048576<9/64$ and
$(1+\theta)^2+\theta^3=66997/32768<9/4$.
$D_0\le37d/64$ and $V/(2m)\le3d/8$ give
$Y\le m^2q+B_0$, $B_0/m\le C=3x/8+109x^2/64$.
For $Q=(2+x-q)\{[m(q+1)+C]^2+2m^2\}$ the same derivative
argument locates the maximum in $[1/2,1]$ with $Q''\le-2m^2$.
Now $c=C/m\le157x/96<1$; the normalized derivative at $3/4$ is
$D=-11/16+(7/2)x+(2x-1)c-c^2$, $|D|\le7/10$.
The lower estimate uses $c+c^2\le(54793/18432)x<(7/2)x$;
the upper uses $c\ge3x/8$ for $x\le2/5$ to get $69/100$, and
$c/x\ge845/672>5/4$ for $x\ge2/5$ to get $43/64$.
The strong-concavity upper polynomial is
\[
\Phi=(5/4+x)\{[7/4+3x/8-3x^2/64]^2+2(1-x^2)^2\}
+(49/400)(1-x^2)^2.
\]
Its ascending coefficients are
\[
\begin{gathered}
10321/1600,\quad429/64,\quad-1903/512-49/200,\\
-4165/1024,\quad40429/16384+49/400,\quad8201/4096.
\end{gathered}
\]
Absorb $x^4\le x^2/4$, $x^5\le x^3/4$ to get
$\Phi\le13/2+(27/4)x-(33/10)x^2-(7/2)x^3$.
This increases to $1056/125<35/4$ on $[0,2/5]$;
afterward use $x^3\ge(2/5)x^2$ to obtain upper bound
$13/2+(27/4)x-(47/10)x^2\le87/10<35/4$.
Thus on the full low rectangle
\begin{equation}\tag{L69.11.7}\label{eq:L69:fivecurvature}
f_2(Y^2+2m^2Y+3m^4)\le(35/4)m^2.
\end{equation}
The full normalized deficit after Taylor expansion is
at least
\[
\begin{gathered}
3173/12800+(591/200)d+(14119/400)d^2+(3307/80)d^3+(423/8)d^4
\\-(21/2)(\sqrt d/m)v-(35/16)(d/m^2)v^2.
\end{gathered}
\]
The same-kernel $\delta_3-\delta_4$ contribution retains the favorable
factor $-(2/5)m^3$ before collecting coefficients.
Using $m^{-1}\le1+4d/3$, $m^{-2}\le16/9$, the first cost has
$x,x^3,x^5,x^7$ coefficients
$2121/1600,8239/800,1057/32,231/8$; the second cost has
$d$ through $d^5$ coefficients
$357035/5760000,\allowbreak 45955/57600,\allowbreak 132265/28800,\allowbreak 15015/1152,\allowbreak 4235/256$.
Use $d^5\le d^4/4$.
On $0\le x\le1/3$, absorbing $x^5\le x^4/3$, $x^7\le x^6/3$
gives a lower polynomial
$3173/12800-(2121/1600)x+(289/100)x^2-(8239/800)x^3+(469/20)x^4$.
Its actual quadratic and quartic coefficients are at least
$16663765/5760000>289/100$ and $1352981/57600>469/20$;
the sixth and eighth remainders are nonnegative.
The quadratic $(469/20)x^2-(8239/800)x+289/100$ is at least $7/4$:
use $8239/800<103/10$ and $(469/5)(57/50)>10609/100$.
The deficit is therefore at least
$3173/12800-(2121/1600)x+(7/4)x^2$, decreasing up to $1/3$,
with endpoint $53/115200>0$.
On $1/3\le x\le\sqrt{1/6}<5/12$, absorb instead using $5/12$.
The lower polynomial is
$Q_*=3173/12800-(2121/1600)x+(289/100)x^2
-(8239/800)x^3+(207/10)x^4$;
the actual quartic coefficient is $1194431/57600>207/10$.
At $1/3$, its value, derivative and second derivative are
$431/345600,1127/4800,5113/400$, all positive, and
$Q_*'''=-24717/400+(2484/5)x>0$ afterward.
Hence it increases and is at least $431/345600>53/115200$.
These overlapping continuous intervals cover $d\le1/6$.
Since $m^4\ge81/256$, their loss proves
$\E[f_2\{T(h_2f_5)\}^4]\le1-d/10000$.
For $d\ge1/6$ the exact actual mean is
\begin{equation*}\tag{L69.11.8}
\E(h_2f_5)=1-4d+2(1+d)\delta_2
\le1-(51/16)d+(13/16)d^2\le283/576<1/2.\label{eq:L69:mean}
\end{equation*}
The quadratic decreases on $[1/6,1/3]$, so $TF\le283/288<1$ and
the reference integral is at most $1-d$.
All stated square inputs and claims use actual fixed directions;
at $d=0$ use original leaf integration, and restore equal pairs and
simultaneous full reversal by original positivity.
